\section{Преговор по комбинаторика}
Тази среща е преговорно занятие върху материал, който се предполага познат от курса по дискретна математика. Основните понятия се припомнят чрез поредица от задачи; част от решенията се извеждат подробно, а останалите са оставени за самостоятелна работа. Сред припомнените инструменти е принципът за включване и изключване, който преподавателят свързва и със задачите от изпита по дискретна математика при доц. Марков.

\subsection{Правило за умножение и пермутации}
Нека разгледаме крайно множество от $n$ елемента. Броят на възможните наредби на тези елемента (пермутации без повторение) се дава от факториела на $n$:
\[ n! = n(n-1) \dots 2 \cdot 1 \]

\begin{supp}[формула на Стирлинг]\label{supp:stirling}
С нарастването на $n$ числата $n!$ растат толкова бързо, че прякото им пресмятане
става практически невъзможно. Формулата на Стирлинг дава удобно приближение:
\[ n! \sim \sqrt{2\pi n}\left(\frac{n}{e}\right)^{n}, \qquad n \to \infty . \]
\end{supp}
По дефиниция приемаме, че $0! = 1$.

\begin{keythm}[Принцип за включване и изключване]\label{thm:incl-excl}
За произволни крайни множества $A_1, \dots, A_n$:
\[
\left| \bigcup_{i=1}^{n} A_i \right|
 = \sum_{i} |A_i|
 - \sum_{i<j} |A_i \cap A_j|
 + \sum_{i<j<k} |A_i \cap A_j \cap A_k|
 - \dots
 + (-1)^{n-1} \left| \bigcap_{i=1}^{n} A_i \right| .
\]
\end{keythm}

За две множества това е познатото $|A \cup B| = |A| + |B| - |A \cap B|$: изваждаме
сечението, защото елементите в него са преброени два пъти. За три множества след като
извадим трите двойни сечения сме извадили тройното сечение един път в повече, затова го
добавяме обратно. Оттук и редуването на знаците.

Доказателството е по индукция: пишем $\bigcup_{i=1}^{n} A_i = \left(\bigcup_{i=1}^{n-1} A_i\right) \cup A_n$,
прилагаме формулата за две множества и след това индукционното предположение — веднъж за
набора $A_1, \dots, A_{n-1}$ и веднъж за набора от сеченията $A_i \cap A_n$. Записът е
тромав, но идеята е тази; важното е да се помни защо знаците се редуват.

\subsection{Вариации без повторение}
Ако от $n$ елемента избираме $k$ и \emph{редът има значение}, но повторения не се допускат,
броят на възможностите е
\[
V_n^k = n(n-1)\dots(n-k+1) = \frac{n!}{(n-k)!} .
\]
Разсъждението е последователно: за първата позиция има $n$ възможности, за втората $n-1$
(един елемент вече е употребен) и така до $k$-тата позиция, за която остават $n-k+1$
възможности. Такива наредени $k$-торки се наричат \emph{вариации} от клас $k$ без повторение.
Забележете, че $V_n^k$ е точно числителят в биномния коефициент по-долу: разликата е
дали редът се брои, или не.

\subsection{Биномни коефициенти (Комбинации)}
\begin{defn}[биномен коефициент]\label{def:binom}
Броят на начините да изберем $k$ елемента от множество с $n$ елемента, като редът на избиране няма значение (комбинации без повторение), се означава с
\[ \binom{n}{k} = C_n^k = \frac{n!}{k!(n-k)!} = \frac{n(n-1)\dots(n-k+1)}{k!} . \]
\end{defn}

Числата от дефиниция~\ref{def:binom} се наричат още биномни коефициенти, тъй като участват в теорема~\ref{thm:newton} за Нютоновия бином.

\begin{prop}[Основни свойства на биномните коефициенти]\label{prop:binom-props}
\begin{enumerate}
    \item \emph{Симетрия:} $\displaystyle \binom{n}{k} = \binom{n}{n-k}$.
    \item \emph{Рекурентна връзка (Правило на Паскал):} $\displaystyle \binom{n}{k} = \binom{n-1}{k-1} + \binom{n-1}{k}$.
    \item \emph{Сума:} $\displaystyle \sum_{k=0}^{n} \binom{n}{k} = 2^n$.
    \item \emph{Знакопроменлива сума:} $\displaystyle \sum_{k=0}^{n} (-1)^k \binom{n}{k} = 0$.
\end{enumerate}
\end{prop}

Свойствата в твърдение~\ref{prop:binom-props}, и в частност правилото на Паскал, позволяват лесно пресмятане на биномните коефициенти и конструирането на Триъгълника на Паскал. Сумата $2^n$ представлява броя на всички подмножества на множество с $n$ елемента; знакопроменливата сума (получена от Нютоновия бином при $a=1$, $b=-1$) означава, че броят на подмножествата с четен брой елементи е равен на броя на подмножествата с нечетен брой елементи.

\subsection{Нютонов бином}

\begin{thm}[Нютонов бином]\label{thm:newton}
\[ (a+b)^n = \sum_{k=0}^{n} \binom{n}{k} a^k b^{n-k} = \binom{n}{0} b^n + \binom{n}{1} ab^{n-1} + \dots + \binom{n}{n} a^n \]
\end{thm}

\subsection{Разпределяне на частици по клетки}
Често в комбинаториката се срещат задачи за разпределяне на частици по клетки.
Ако имаме $k$ различими частици и $n$ клетки, броят на възможните разпределения е:
\[ n^k \]
Това е така, защото всяка от $k$-те частици може да бъде поставена в коя да е от $n$-те клетки напълно независимо от останалите.

\begin{example}[Разпределяне без празна клетка]\label{ex:15-surj}
Нека $k \ge n$ и разпределяме $k$ различими частици в $n$ различими клетки така, че
\emph{никоя клетка да не остане празна}. Броят на такива разпределения е
\[
\sum_{m=0}^{n} (-1)^m \binom{n}{m} (n-m)^k .
\]

Пряко преброяване не работи: условието „няма празна клетка“ не се разлага на независими
избори. Затова броим допълнението чрез принцип~\ref{thm:incl-excl}. Нека
$A_i$ е множеството от разпределения, при които $i$-тата клетка е празна. Тогава
разпределенията с поне една празна клетка са $\left| \bigcup_{i=1}^n A_i \right|$.

Ако фиксираме $m$ конкретни клетки и искаме всички те да са празни, остават $n-m$ клетки за
$k$-те частици, тоест $(n-m)^k$ разпределения; а $m$-те клетки се избират по
$\binom{n}{m}$ начина. По принципа за включване и изключване
\[
\left| \bigcup_{i=1}^n A_i \right| = \sum_{m=1}^{n} (-1)^{m-1} \binom{n}{m} (n-m)^k ,
\]
и като извадим това от общия брой $n^k$, получаваме твърдението (членът $m=0$ дава точно
$n^k$).
\end{example}

Този пример е добра илюстрация защо изобщо съществува принципът за включване и изключване:
когато налагаме условие от вида „нищо не е празно“, отделните случаи се припокриват и
трябва да се коригират последователно.

\subsection{Уравнения с цели неотрицателни числа (Метод „Звезди и прегради“)}
Разглеждаме уравнението:
\[ x_1 + x_2 + \dots + x_k = n \]
където $x_i$ са цели неотрицателни числа ($x_i \geq 0$ за всяко $i=1,\dots,k$) и $n \geq 0$ е цяло число.
Търсим броя на решенията на това уравнение. Можем да си представим, че имаме $n$ неразличими обекта (единици или „звезди“), които искаме да разпределим в $k$ клетки. За да разделим обектите на $k$ групи, ни трябват $k-1$ „прегради“ (или „черти“).

Общият брой на позициите за звездите и преградите е $n + k - 1$. От тези позиции трябва да изберем $k-1$ места, на които да поставим преградите (или еквивалентно, да изберем $n$ места за звездите).

\begin{prop}[Звезди и прегради]\label{prop:stars-bars}
Броят на решенията на уравнението $x_1 + x_2 + \dots + x_k = n$ в цели неотрицателни числа е
\[ \binom{n+k-1}{k-1} = \binom{n+k-1}{n} . \]
\end{prop}

\begin{cor}[Комбинации с повторение]\label{cor:multiset}
Броят на начините да изберем $k$ елемента от $n$-елементно множество, когато редът няма
значение, но \emph{повторенията са позволени} (тоест броят на $k$-елементните
мултимножества), е
\[ \binom{n+k-1}{k} . \]
\end{cor}

\begin{proof}
Да изберем мултимножество означава да кажем колко пъти участва всеки от $n$-те елемента:
това е избор на цели неотрицателни числа $x_1, \dots, x_n$ с $x_1 + \dots + x_n = k$.
По твърдение~\ref{prop:stars-bars} с $k$ на брой звезди и $n$ клетки броят е
$\binom{k+n-1}{n-1} = \binom{n+k-1}{k}$.
\end{proof}

Сравнете двете съседни формули: $\binom{n}{k}$ брои подмножествата (без повторение), а
$\binom{n+k-1}{k}$ — мултимножествата (с повторение).

Твърдение~\ref{prop:stars-bars} е изключително важно и се прилага в много контексти, включително за намиране на броя на разпределенията на $n$ неразличими частици в $k$ различими клетки (където редът на частиците в дадена клетка няма значение).

\section{Просто случайно блуждаене}

Ще разгледаме един класически модел в теорията на вероятностите — просто случайно блуждаене (Simple Random Walk).
Нека една частица се движи по целочислените точки на реалната права. Тя стартира от началото на координатната система (точка 0) в момент $t=0$. Във всеки следващ дискретен момент от времето $t=1, 2, \dots$ частицата прави стъпка надясно (към $+1$) или наляво (към $-1$).
Нека $\xi_i \in \{+1, -1\}$ е стойността на $i$-тата стъпка. Тогава позицията на частицата след $n$ стъпки е:
\[ S_n = \xi_1 + \xi_2 + \dots + \xi_n \]
като по дефиниция приемаме $S_0 = 0$.

Можем да представим блуждаенето графично в координатната равнина $(t, S_t)$, където по хоризонталната ос нанасяме времето (броя стъпки), а по вертикалната — позицията на частицата. Тогава блуждаенето представлява начупена линия, започваща от $(0,0)$. Всяка отсечка от тази линия свързва точките $(i-1, S_{i-1})$ и $(i, S_i)$.

\subsection{Брой на всички пътища}
Ако разгледаме блуждаене с дължина $n$, то броят на всички възможни пътища (траектории) е $2^n$, тъй като на всяка от $n$-те стъпки имаме точно 2 възможни избора ($+1$ или $-1$).

За да стигнем от точка $(0,0)$ до точка $(n, x)$, трябва броят на стъпките надясно ($+1$) и наляво ($-1$) да са съответно $p$ и $q$, такива че да са изпълнени следните две условия:
\begin{align*}
p + q &= n \\
p - q &= x
\end{align*}
Като съберем двете уравнения, намираме:
\[ 2p = n + x \implies p = \frac{n+x}{2} \]
За да съществува такъв път, е необходимо и достатъчно $n+x$ да бъде четно число и $-n \leq x \leq n$.
\begin{prop}[Брой пътища до дадена точка]\label{prop:paths-count}
Ако $n+x$ е четно и $-n \le x \le n$, броят на пътищата от $(0,0)$ до $(n,x)$ е равен на броя на начините да изберем $p = \frac{n+x}{2}$ стъпки надясно от общо $n$ налични стъпки:
\[ N_n(x) = \binom{n}{p} = \binom{n}{\frac{n+x}{2}} . \]
\end{prop}

\section{Принцип на отражението}
Често ни интересуват пътища, които удовлетворяват определени условия, например пътища, които остават неотрицателни ($S_i \geq 0$ за всяко $i$). За да преброим такива пътища, използваме Принципа на отражението, въведен от Дезире Андре.

\begin{keythm}[Принцип на отражението]\label{thm:reflection}
Нека $x, y > 0$. Броят на пътищата от $(0,x)$ до $(n,y)$, които докосват или пресичат хоризонталната ос (т.е. $S_i = 0$ за някое $i$), е равен на общия брой пътища от $(0,-x)$ до $(n,y)$.
\end{keythm}

\begin{lecfigure}[Принципът на отражението. Плътната линия е „лош“ път от $(0,x)$ до
$(n,y)$; $t_0$ е първият момент, в който той докосва оста. Отразявайки началната част до момента $t_0$ спрямо оста (прекъснатата линия), получаваме път от $(0,-x)$ до $(n,y)$.
Съответствието е биекция.\label{fig:reflection}]
\begin{tikzpicture}[x=0.62cm, y=0.52cm]
  % axes
  \draw[axisline] (-0.4,0) -- (9.4,0) node[right, font=\small] {$t$};
  \draw[axisline] (0,-3.6) -- (0,3.9) node[above, font=\small] {$S$};
  % the "bad" path: (0,2) -> touches 0 at t=4 -> (8,2)
  \draw[walk] (0,2) -- (1,3) -- (2,2) -- (3,1) -- (4,0) -- (5,1) -- (6,2) -- (7,1) -- (8,2);
  % reflection of the initial segment across S = 0
  \draw[walkrefl] (0,-2) -- (1,-3) -- (2,-2) -- (3,-1) -- (4,0);
  % guides
  \draw[guide] (4,0) -- (4,-3.4);
  \node[font=\small] at (4,-3.9) {$t_0$};
  % endpoints
  \fill[thmframe] (0,2) circle (2.1pt) node[left, font=\small, black] {$x$};
  \fill[thmframe] (8,2) circle (2.1pt) node[right, font=\small, black] {$y$};
  \fill[keyframe] (0,-2) circle (2.1pt) node[left, font=\small, black] {$-x$};
  \fill[gray!70] (4,0) circle (1.8pt);
\end{tikzpicture}
\end{lecfigure}

Нека наречем пътищата, които докосват оста, „лоши“ (фигура~\ref{fig:reflection}). Вземаме един произволен такъв лош път. Нека $t_0$ е първият момент, в който пътят докосва нулата (т.е. $S_{t_0} = 0$, но $S_i > 0$ за $i < t_0$). Принципът на отражението гласи, че ако отразим симетрично началната част от пътя (от $t=0$ до $t=t_0$) спрямо абсцисната ос $S=0$, ще получим нов път, който започва от $(0,-x)$, докосва нулата в $t_0$ и продължава по оригиналния път до $(n,y)$.

Съществено е, че отразяваме спрямо \emph{първия} момент на докосване. Ако избирахме
произволен момент, в който пътят е на нулата, един и същ лош път би давал по няколко
различни образа и съответствието нямаше да е еднозначно. Първият момент е определен
еднозначно от самия път, затова изображението е коректно дефинирано; обратно, при даден
път от $(0,-x)$ до $(n,y)$ неговият пръв момент на докосване възстановява еднозначно
оригинала.

Тъй като всеки път от $(0,-x)$ до $(n,y)$ (при $x,y > 0$) задължително пресича оста $S=0$ в някакъв момент (започва под нея и завършва над нея), съществува взаимно еднозначно съответствие (биекция) между лошите пътища от $(0,x)$ до $(n,y)$ и множеството на всички възможни пътища от $(0,-x)$ до $(n,y)$. Това доказва твърдението.

\begin{supp}[Отражение спрямо произволно ниво]\label{supp:reflection-level}
Теорема~\ref{thm:reflection} е формулирана за ниво $0$ и изисква $x, y > 0$. Същият
аргумент работи спрямо кое да е ниво $a$, стига началната и крайната точка да лежат от
една и съща страна на него: ако $x, y > a$, то броят на пътищата от $(0,x)$ до $(n,y)$,
които докосват ниво $a$, е равен на общия брой пътища от $(0,\, 2a - x)$ до $(n,y)$.
Наистина, $2a-x$ е огледалният образ на $x$ спрямо $a$, и отразяването на началната част
до първия момент на докосване на ниво $a$ дава същата биекция.

Тази форма ни е нужна веднага по-долу: там началната и крайната точка са на самата ос
($x = y = 0$), така че теорема~\ref{thm:reflection} не се прилага пряко, а работим с
ниво $a = -1$.
\end{supp}

\subsection{Пътища, които остават неотрицателни}
Нека разгледаме пътищата, стартиращи от $(0,0)$ и завършващи в $(2n, 0)$. Според предходното твърдение общият им брой е $N_{2n}(0) = \binom{2n}{n}$.

Интересуваме се от броя на пътищата от $(0,0)$ до $(2n, 0)$, които остават неотрицателни през цялото време, т.е. $S_i \geq 0$ за $i=1, \dots, 2n$.
За да намерим техния брой, ще преброим „лошите“ пътища — тези, за които $S_i = -1$ за поне едно $i$ — и ще ги извадим от общия брой пътища.

Тук началната и крайната точка лежат на самата ос, затова теорема~\ref{thm:reflection}
не може да се приложи буквално. Използваме вариантa от допълнение~\ref{supp:reflection-level}
с ниво $a = -1$: началото $0$ и краят $0$ са от една и съща страна на $-1$, а „лоши“ са
точно пътищата, които докосват ниво $-1$. Огледалният образ на началната точка спрямо
$-1$ е $2(-1) - 0 = -2$, следователно лошите пътища са в биекция с \emph{всички} пътища
от $(0,-2)$ до $(2n,0)$.

Броят на пътищата от $(0,-2)$ до $(2n,0)$ се намира като решим:
\begin{align*}
p + q &= 2n \\
p - q &= 2
\end{align*}
Оттук $2p = 2n + 2 \implies p = n+1$. Следователно броят им е $\binom{2n}{n+1}$.

Тогава броят на „добрите“ пътища (които никога не достигат $-1$, следователно $S_i \geq 0$ за всички $i$) от $(0,0)$ до $(2n,0)$ е:
\[ \binom{2n}{n} - \binom{2n}{n+1} \]
Тази разлика може да се опрости алгебрично по следния начин:
\begin{align*}
\binom{2n}{n} - \binom{2n}{n+1} &= \frac{(2n)!}{n!n!} - \frac{(2n)!}{(n+1)!(n-1)!} \\
&= \frac{(2n)!}{n!(n-1)!} \left( \frac{1}{n} - \frac{1}{n+1} \right) \\
&= \frac{(2n)!}{n!(n-1)!} \left( \frac{n+1 - n}{n(n+1)} \right) \\
&= \frac{(2n)!}{n!(n-1)!} \frac{1}{n(n+1)} \\
&= \frac{1}{n+1} \frac{(2n)!}{n!n!} \\
&= \frac{1}{n+1} \binom{2n}{n}
\end{align*}
Този забележителен резултат представлява $n$-тото число на Каталан ($C_n$). Числата на Каталан имат огромно значение в комбинаториката, тъй като те броят множество различни комбинаторни структури (например правилни скобови изрази, триангулации на изпъкнали многоъгълници и дървета).

\subsection{Пътища с условие за строга положителност}
Ако разгледаме по-строго условие — да намерим броя на пътищата от $(0,0)$ до $(2n,0)$, които се връщат в нулата за първи път чак на последната стъпка $2n$ (т.е. $S_i > 0$ за $1 \leq i \leq 2n-1$ и $S_{2n} = 0$).

Тогава първата стъпка на такъв път задължително е нагоре към $(1,1)$, а предпоследната позиция задължително е $(2n-1, 1)$, откъдето прави стъпка надолу към $(2n,0)$.
Броят на тези пътища съвпада с броя на пътищата от $(1,1)$ до $(2n-1, 1)$, които никога не докосват нулата (остават строго положителни, $S_i > 0$).

Ако транслираме координатната система с вектор $(-1,-1)$, това е равносилно на намирането на броя на пътищата от $(0,0)$ до $(2n-2, 0)$, които остават неотрицателни (т.е. не слизат под новото ниво 0).
Според предходната точка, този брой е равен на съответното число на Каталан за $2n-2$ стъпки:
\[ \frac{1}{n} \binom{2n-2}{n-1} = C_{n-1} \]
което е $(n-1)$-тото число на Каталан.

\section{Задачи}

Следващите условия възстановяват задачите, които са формулирани устно или са записани на дъската. Запазена е номерацията от занятието.

Условията на задачи от 6 до 9 и 11 не са произнесени или записани в наличните материали: преподавателят се позовава на предварително раздаден лист със задачи и разработва само задача 10. Затова липсващите условия не могат да бъдат възстановени достоверно.

\begin{enumerate}
    \item[1.] Припомнете принципа за включване и изключване за $n$ крайни множества. Доказателството е по индукция, като индукционното предположение трябва да се приложи и към подходящото семейство от сечения; за този курс е важно преди всичко да се помни кога и как се използва принципът.

    \item[2.] Нека $M$ е множество с $n$ елемента. Намерете броя на:
    \begin{enumerate}
        \item $k$-елементните подмножества на $M$;
        \item наредените $k$-торки от различни елементи на $M$;
        \item наредените $k$-торки от елементи на $M$, когато повторенията са позволени.
    \end{enumerate}

    \item[3.] Намерете броя на решенията на
    \[
    x_1+\dots+x_k=n
    \]
    първо при $x_1,\dots,x_k\in\mathbb{N}$, а след това при $x_1,\dots,x_k\in\mathbb{N}_0$. Използвайте метода „звезди и прегради“.

    \item[4.] Разпределяме $k$ различими частици в $n$ различими клетки. Намерете броя на разпределенията, когато:
    \begin{enumerate}
        \item всяка клетка съдържа най-много една частица;
        \item няма ограничение за броя частици в клетка;
        \item няма празна клетка, като $k\ge n$. За последната точка използвайте принципа за включване и изключване.
    \end{enumerate}

    \item[5.] Разпределяме $k$ неразличими частици в $n$ различими клетки. Намерете броя на разпределенията, когато:
    \begin{enumerate}
        \item всяка клетка съдържа най-много една частица;
        \item няма ограничение за броя частици в клетка;
        \item няма празна клетка. Последната точка е оставена за домашна работа.
    \end{enumerate}

    \item[10.] Разгледайте просто случайно блуждаене, което започва от $(0,0)$ и на всяка стъпка се изменя с $+1$ или $-1$. Пребройте всички пътища до $(2n,0)$, а след това пътищата, които остават неотрицателни. За второто преброяване приложете принципа на отражението.
\end{enumerate}
